\chapter{Αποτελέσματα}
\label{ch:results}

\section{Στατικό σενάριο αναφοράς}
\label{sec:results-static}

Ο Πίνακας~\ref{tab:static-scenario} παρουσιάζει την έξοδο μιας μεμονωμένης
εκτέλεσης του σεναρίου αναφοράς (\texttt{rng(42)}), με τον κανόνα
πολυσυνδεσιμότητας της εξίσωσης~\eqref{eq:node-selection}: οι τέσσερις
χρήστες κοντά στους δύο σταθμούς βάσης εξυπηρετούνται πλέον
\emph{ταυτόχρονα} από τον καλύτερο BS \emph{και} τον δορυφόρο
(DualConnectivity), ενώ οι δύο απομακρυσμένοι χρήστες (100+~km από κάθε
σταθμό βάσης) εξυπηρετούνται μόνο από τον δορυφόρο, με γωνία ανύψωσης
$78^\circ$-$89^\circ$ (σχεδόν στο ζενίθ) και για τους έξι.

\begin{table}[htbp]
\centering
\caption{Αποτελέσματα στατικού σεναρίου αναφοράς (μία εκτέλεση, \texttt{rng(42)}), με πολυσυνδεσιμότητα. Cand.~BS/Sat.~SNR: SNR του καλύτερου υποψήφιου BS / δορυφόρου, ανεξάρτητα από το αν χρησιμοποιείται τελικά.}
\label{tab:static-scenario}
\small
\begin{tabular}{@{}c c c c c c c@{}}
\toprule
User & Node & Type & Cand.~BS SNR (dB) & Cand.~Sat.~SNR (dB) & Capacity (Mbps) & Energy (\textmu J/bit) \\
\midrule
1 & BS1+SAT-1 & DualConnectivity & 33.83  & 6.61 & 110.25 & 2.0392 \\
2 & BS2+SAT-1 & DualConnectivity & 57.24  & 6.61 & 42.26  & 1.7897 \\
3 & BS2+SAT-1 & DualConnectivity & 39.67  & 6.61 & 42.26  & 1.7897 \\
4 & BS2+SAT-1 & DualConnectivity & 5.36   & 6.61 & 21.42  & 3.5306 \\
5 & SAT-1     & Satellite        & $-612.1$ & 6.43 & 8.10 & 0.1292 \\
6 & SAT-1     & Satellite        & $-58.68$ & 6.61 & 8.27 & 0.1266 \\
\bottomrule
\end{tabular}
\end{table}

Ο χρήστης~$4$ είναι το πιο διαφωτιστικό παράδειγμα της αλλαγής κανόνα: το
BS SNR του ($5.36$~dB) είναι \emph{χαμηλότερο} από το SNR του δορυφόρου
($6.61$~dB) -- κάτω από τον παλιό, single-connectivity κανόνα $\argmax$
θα είχε ανατεθεί \emph{μόνο} στον δορυφόρο, όπως συνέβαινε πριν την
εισαγωγή της πολυσυνδεσιμότητας. Και τα δύο SNR ξεπερνούν όμως το
ελάχιστο χρησιμοποιήσιμο κατώφλι $\text{SNR}_{\min} \approx -7.53$~dB,
οπότε ο νέος κανόνας τον συνδέει και στους δύο ταυτόχρονα: η
χωρητικότητά του ($21.42$~Mbps) είναι το άθροισμα μιας μικρής επίγειας
συνεισφοράς (μοιρασμένη σε $3$ χρήστες στο BS2, με χαμηλή φασματική
απόδοση λόγω του οριακού SNR) και της δορυφορικής συνεισφοράς, ενώ η
ενέργειά του ανά bit ($3.5306$~\textmu J/bit) είναι η υψηλότερη μεταξύ
των DualConnectivity χρηστών -- πληρώνει την ισχύ \emph{δύο} ενεργών
ραδιοσυστημάτων για μια σχετικά μικρή αύξηση χωρητικότητας. Αντίθετα, οι
χρήστες $5$-$6$, με βαθιά αρνητικό BS SNR (εκτός εμβέλειας/NLOS σε
$100+$~km), παραμένουν Satellite-only: το κατώφλι τούς αποκλείει από τη
DualConnectivity, όχι επειδή δεν προσπάθησαν, αλλά επειδή ο επίγειος
υποψήφιός τους είναι πραγματικά μη-χρησιμοποιήσιμος (βλ.
Ενότητα~\ref{subsec:results-outage} για το πλήρες σκεπτικό του
κατωφλίου).

Οι χρήστες $2$-$3$, που μοιράζονται το BS2, καταλήγουν στην \emph{ίδια}
επίγεια συνεισφορά παρότι το BS SNR τους διαφέρει κατά $17.6$~dB -- και οι
δύο είναι κορεσμένοι στο όριο MCS/CQI (εξίσωση~\eqref{eq:shannon-capacity},
$\eta_{\max} = 5.5547$~bit/s/Hz), οπότε η χωρητικότητά τους καθορίζεται
αποκλειστικά από το μερίδιο εύρους ζώνης, όχι πια από το SNR· ο
χρήστης~$1$, μόνος στο $B_{\text{BS}}$ του BS1, φτάνει την υψηλότερη
συνολική χωρητικότητα ($110.25$~Mbps).

Το σενάριο αυτό είναι χρήσιμο ως σημείο αναφοράς και ως βάση ελέγχου
ορθότητας της υλοποίησης, παραμένει όμως μια \emph{μεμονωμένη} στοχαστική
πραγματοποίηση: μια διαφορετική επιλογή LOS/NLOS ή shadow fading σε
οποιαδήποτε ζεύξη θα έδινε διαφορετικούς αριθμούς. Οι Ενότητες
\ref{sec:results-repeated} και \ref{sec:results-temporal} επεκτείνουν την
ανάλυση πέρα από αυτή τη μεμονωμένη στιγμιότυπη εικόνα.

\section{Στατιστική ανάλυση KPI (500 επαναλαμβανόμενα runs)}
\label{sec:results-repeated}

Για να χαρακτηριστεί η διακύμανση των KPI που οφείλεται αποκλειστικά στη
στοχαστικότητα του καναλιού (LOS draw και shadow fading, ενότητα
\ref{sec:meth-terrestrial}), η \emph{ίδια} τοπολογία εκτελέστηκε $500$
φορές με διαφορετικό RNG seed ανά φορά (\texttt{kpiRepeatedRuns.m}). Με
τον κανόνα πολυσυνδεσιμότητας (εξίσωση~\eqref{eq:node-selection}), οι
τέσσερις κοντινοί χρήστες καταλήγουν \texttt{DualConnectivity} σε
\emph{κάθε} μία από τις $500$ επαναλήψεις -- το BS SNR τους παραμένει
πάντα αρκετά πάνω από $\text{SNR}_{\min}$ ώστε η στοχαστικότητα του
καναλιού να μην το ρίξει ποτέ κάτω από το κατώφλι -- οπότε η κατηγορία
"Terrestrial-only" δεν εμφανίζεται καθόλου σε αυτό το σύνολο
επαναλήψεων· ο Πίνακας~\ref{tab:repeated-runs-summary} συνοψίζει
Throughput/Energy ανά τις δύο κατηγορίες που όντως εμφανίζονται.
Το SNR δεν έχει νόημα ως ένα ενιαίο νούμερο για μια DualConnectivity
γραμμή (δύο ζεύξεις πολύ διαφορετικής κλίμακας ταυτόχρονα), οπότε
αναφέρεται ξεχωριστά, ανά \emph{ζεύξη} αντί ανά νικητή, στον
Πίνακα~\ref{tab:link-snr-std}.

\begin{table}[htbp]
\centering
\caption{Στατιστική σύνοψη KPI πάνω από 500 επαναλήψεις της σταθερής
τοπολογίας αναφοράς (mean $\pm$ std, [min, max]). Outage: $0/3000$ γραμμές.}
\label{tab:repeated-runs-summary}
\small
\begin{tabular}{@{}l c c@{}}
\toprule
KPI & DualConnectivity ($n=1991$) & Satellite ($n=1009$) \\
\midrule
Throughput (Mbps)       & $52.0 \pm 23.5$ \, [9.9, 110.3]   & $8.19 \pm 0.08$ \, [8.10, 8.27] \\
Energy per bit (\textmu J/bit) & $2.61 \pm 1.59$ \, [1.69, 12.06]    & $0.1279 \pm 0.0013$ \, [0.1266, 0.1292] \\
\bottomrule
\end{tabular}
\end{table}

\begin{table}[htbp]
\centering
\caption{Τυπική απόκλιση SNR ανά \emph{ζεύξη} (όχι ανά νικητή κόμβο),
πάνω στους χρησιμοποιήσιμους υποψηφίους των ίδιων 500 επαναλήψεων --
πηγή του noise sigma της Ενότητας~\ref{sec:results-ml-realistic}.}
\label{tab:link-snr-std}
\begin{tabular}{@{}l c@{}}
\toprule
Ζεύξη & std SNR (dB) \\
\midrule
Terrestrial (καλύτερο BS)  & $19.67$ \\
Satellite                   & $0.0682$ \\
\bottomrule
\end{tabular}
\end{table}

\begin{figure}[htbp]
\centering
\includegraphics[width=0.75\textwidth]{kpi_hist_Capacity_Mbps.png}
\caption{Κατανομή throughput (Mbps) ανά τύπο εξυπηρέτησης, πάνω από 500 επαναλήψεις.}
\label{fig:kpi-hist-capacity}
\end{figure}

\begin{figure}[htbp]
\centering
\includegraphics[width=0.75\textwidth]{kpi_hist_EnergyPerBit_uJ.png}
\caption{Κατανομή ενέργειας ανά bit (\textmu J/bit) ανά τύπο εξυπηρέτησης, πάνω από 500 επαναλήψεις.}
\label{fig:kpi-hist-energy}
\end{figure}

\begin{figure}[htbp]
\centering
\includegraphics[width=0.75\textwidth]{kpi_hist_SNR_dB.png}
\caption{Κατανομή SNR (dB) του νικητή κόμβου, πάνω από 500 επαναλήψεις.
Μόνο η κατηγορία Satellite εμφανίζει ουσιαστική κατανομή εδώ -- ένας
DualConnectivity χρήστης δεν έχει ένα ενιαίο "SNR νικητή" (η στήλη
\texttt{SNR\_dB} είναι \texttt{NaN} για αυτές τις γραμμές, βλ.
Ενότητα~\ref{sec:meth-node-selection}), και η κατηγορία Terrestrial-only
δεν εμφανίζεται καθόλου (Πίνακας~\ref{tab:repeated-runs-summary}). Η
διακύμανση SNR ανά ζεύξη (αντί ανά νικητή) δίνεται στον
Πίνακα~\ref{tab:link-snr-std}.}
\label{fig:kpi-hist-snr}
\end{figure}

Δύο παρατηρήσεις προκύπτουν άμεσα:

\begin{itemize}
  \item \textbf{Ασύμμετρη ενεργειακή απόδοση.} Ένας Satellite-only
    χρήστης έχει $\sim\!6.3\times$ χαμηλότερο μέσο throughput από έναν
    DualConnectivity χρήστη, αλλά $\sim\!20\times$ χαμηλότερη ενέργεια
    ανά bit. Αυτό προκύπτει άμεσα από την εξίσωση~\eqref{eq:energy-per-bit}:
    η πολύ μικρότερη κατανάλωση ισχύος υποδομής του δορυφορικού payload
    (Ενότητα~\ref{sec:meth-energy}, $P_{\text{sat}} \approx 6.3$~W
    συνολικά, μοιρασμένα σε όλους τους συνδεδεμένους στον δορυφόρο --
    Satellite-only \emph{και} DualConnectivity χρήστες μαζί -- έναντι
    $P_{\text{BS}} \approx 223.8$~W ανά ενεργό BS) υπερκαλύπτει τον
    χαμηλότερο ρυθμό μετάδοσης. Ένας DualConnectivity χρήστης, αντίθετα,
    πληρώνει το ενεργειακό κόστος \emph{και} των δύο ζεύξεών του
    (εξίσωση~\eqref{eq:energy-per-bit}, άθροισμα πάνω σε
    $\mathcal{N}(u)$) για μια χωρητικότητα που κυριαρχείται σχεδόν
    αποκλειστικά από την επίγεια συνεισφορά του -- η δορυφορική ζεύξη
    προσθέτει μόλις μερικά Mbps πάνω σε δεκάδες, αλλά προσθέτει πλήρως
    το δικό της ενεργειακό μερίδιο, γι' αυτό και η ενέργεια ανά bit ενός
    DualConnectivity χρήστη ($2.61$~\textmu J/bit μέσος όρος) δεν
    διαφέρει πολύ από αυτή ενός καθαρά επίγειου χρήστη πριν την
    πολυσυνδεσιμότητα.
  \item \textbf{Ασύμμετρη διακύμανση SNR ανά ζεύξη.} Ο
    Πίνακας~\ref{tab:link-snr-std} δείχνει ότι η τυπική απόκλιση SNR στο
    επίγειο σκέλος ($19.67$~dB) είναι περίπου δύο τάξεις μεγέθους
    μεγαλύτερη από αυτή του δορυφορικού σκέλους ($0.068$~dB). Πρόκειται
    για άμεση συνέπεια της μοντελοποίησης: το επίγειο κανάλι διαθέτει
    στοχαστικό LOS/NLOS draw και shadow fading
    (εξισώσεις~\eqref{eq:los-prob-uma}-\eqref{eq:pathloss-shadowing}),
    ενώ το δορυφορικό μοντέλο (εξισώσεις~\eqref{eq:fspl}-\eqref{eq:gas-atten})
    παραμένει καθαρά ντετερμινιστικό δεδομένης της γεωμετρίας -- η
    προσθήκη της ατμοσφαιρικής απόσβεσης αερίων δεν εισάγει καμία
    στοχαστικότητα, μόνο μια μικρή, ντετερμινιστική διόρθωση. Ο δορυφόρος
    εμφανίζεται με άλλα λόγια τεχνητά "σταθερός", επειδή το μοντέλο δεν
    του αποδίδει καμία στοχαστική συνιστώσα -- ένα όριο της τρέχουσας
    μοντελοποίησης που συζητείται στο Κεφάλαιο~\ref{ch:discussion}. Αυτή
    η ίδια, μετρημένη διακύμανση είναι η πηγή του θορύβου SNR στο τρίτο
    πείραμα μηχανικής μάθησης (Ενότητα~\ref{sec:results-ml-realistic}).
\end{itemize}

\section{Χρονική εξέλιξη διέλευσης LEO}
\label{sec:results-temporal}

Η χρονική προσομοίωση (\texttt{temporalPassSimulation.m}, Ενότητα
\ref{sec:meth-temporal}) εκτέλεσε $93$ διαδοχικά χρονικά βήματα
($\Delta t = 5$~s, συνολική διάρκεια $\approx 460$~s), με το υποδορυφορικό
σημείο να κινείται με ταχύτητα $\approx 6986$~m/s (τροχιακή περίοδος
$\approx 5730$~s). Το Σχήμα~\ref{fig:temporal-elevation} δείχνει το
προκύπτον προφίλ γωνίας ανύψωσης: μια καθαρή, αναμενόμενη καμπύλη διέλευσης,
με μέγιστη γωνία ανύψωσης $\approx 89.9^\circ$ (σχεδόν ζενίθ) γύρω στα
$t \approx 215$~s.

\begin{figure}[htbp]
\centering
\includegraphics[width=0.8\textwidth]{temporal_SatElevation_deg.png}
\caption{Γωνία ανύψωσης δορυφόρου ανά χρήστη κατά τη διάρκεια της διέλευσης (93 βήματα, $\Delta t = 5$~s).}
\label{fig:temporal-elevation}
\end{figure}

\begin{figure}[htbp]
\centering
\includegraphics[width=0.8\textwidth]{temporal_CandBS_SNR_dB.png}
\caption{SNR του καλύτερου υποψήφιου BS ανά χρήστη κατά τη διάρκεια της
διέλευσης (ανεξάρτητα από το αν χρησιμοποιείται τελικά -- σταθερό, αφού
ούτε οι χρήστες ούτε οι σταθμοί βάσης μετακινούνται). Οι χρήστες $1$-$4$
έχουν σταθερά χρησιμοποιήσιμο BS SNR καθ' όλη τη διέλευση· οι χρήστες
$5$-$6$ έχουν βαθιά αρνητικό BS SNR (εκτός εμβέλειας/NLOS), πάντα κάτω
από $\text{SNR}_{\min}$.}
\label{fig:temporal-bs-snr}
\end{figure}

\begin{figure}[htbp]
\centering
\includegraphics[width=0.8\textwidth]{temporal_CandSat_SNR_dB.png}
\caption{SNR του δορυφόρου ανά χρήστη κατά τη διάρκεια της διέλευσης
(σχεδόν ταυτόσημο και για τους $6$ χρήστες, αφού βρίσκονται όλοι κοντά
μεταξύ τους σε σχέση με το υψόμετρο του δορυφόρου). Η γραμμή διακόπτεται
όποτε ο δορυφόρος είναι εκτός ορατότητας
($\text{SNR}_{\text{sat}} = -\infty$). Η μετάβαση κάτω από
$\text{SNR}_{\min} \approx -7.53$~dB προς το τέλος της διέλευσης οδηγεί
τους χρήστες $1$-$4$ σε SN event (αποβολή του δορυφόρου, παραμονή στο
BS) και τους χρήστες $5$-$6$ σε ρητή κατάσταση outage
(Ενότητα~\ref{subsec:results-outage}).}
\label{fig:temporal-sat-snr}
\end{figure}

\subsection{Handovers, SN events και outage events}
\label{subsec:results-handovers}

Ο Πίνακας~\ref{tab:handovers} καταγράφει τρεις ξεχωριστούς τύπους
μετάβασης ανά χρήστη κατά τη διάρκεια της διέλευσης: \emph{handover}
(αλλάζει ο σταθμός βάσης -- ο "master node" MN, στην ορολογία MR-DC του
Majamaa~\cite{majamaa2024toward} που ήδη αναφέρθηκε στο
Κεφάλαιο~\ref{ch:background}), \emph{SN event} (προστίθεται ή αφαιρείται
ο δορυφόρος -- ο "secondary node" SN -- ενώ ο MN παραμένει ο ίδιος), και
\emph{outage event} (μετάβαση προς \texttt{ServingNode="None"}). Η
διάκριση αυτή είναι απαραίτητη μετά την πολυσυνδεσιμότητα
(εξίσωση~\eqref{eq:node-selection}): μια μετάβαση
DualConnectivity$\rightarrow$Terrestrial δεν αλλάζει τον σταθμό βάσης
που εξυπηρετεί τον χρήστη, οπότε δεν είναι πραγματικό handover, παρότι η
τιμή του \texttt{ServingType} αλλάζει.

\begin{table}[htbp]
\centering
\caption{Αριθμός handovers, SN events και outage events ανά χρήστη κατά
τη διάρκεια της διέλευσης (93 διαδοχικά βήματα), με το χωρικά
συσχετισμένο μοντέλο shadow fading (εξίσωση~\eqref{eq:correlated-sf}),
πολυσυνδεσιμότητα και ρητή κατάσταση outage
(εξίσωση~\eqref{eq:node-selection}).}
\label{tab:handovers}
\begin{tabular}{@{}c c c c c c@{}}
\toprule
User & Handovers & SN events & Outage events & Αρχικός κόμβος & Τελικός κόμβος \\
\midrule
1 & 0 & 1 & 0 & BS2+SAT-1 & BS2 \\
2 & 0 & 1 & 0 & BS2+SAT-1 & BS2 \\
3 & 0 & 1 & 0 & BS1+SAT-1 & BS1 \\
4 & 0 & 1 & 0 & BS2+SAT-1 & BS2 \\
5 & 0 & 0 & 1 & SAT-1     & None (outage) \\
6 & 0 & 0 & 1 & SAT-1     & None (outage) \\
\bottomrule
\end{tabular}
\end{table}

Η εικόνα εδώ διαφέρει σημαντικά από αυτή που έδινε τόσο το αρχικό,
ανεξάρτητο-ανά-βήμα μοντέλο σκίασης, όσο και το προγενέστερο,
single-connectivity μοντέλο επιλογής κόμβου. Δεν καταγράφεται \emph{κανένα}
πραγματικό handover: οι χρήστες $1$-$4$ (κοντά σε σταθμό βάσης καθ' όλη τη
διέλευση) διατηρούν σταθερά τον ίδιο "master" σταθμό βάσης, κάτι
αναμενόμενο αφού ούτε οι ίδιοι ούτε οι σταθμοί βάσης μετακινούνται· δεν
υπάρχει δηλαδή φυσικός λόγος να αλλάζει ο νικητής BS1/BS2 ανά βήμα, και το
συσχετισμένο μοντέλο πλέον το αντικατοπτρίζει σωστά. Αντ' αυτού, οι
χρήστες $1$-$4$ ξεκινούν \emph{ήδη} σε DualConnectivity (το BS SNR τους
είναι σταθερά αρκετά καλό, Σχήμα~\ref{fig:temporal-bs-snr}, ενώ ο
δορυφόρος είναι ήδη ορατός/χρησιμοποιήσιμος από την αρχή της διέλευσης,
Σχήμα~\ref{fig:temporal-sat-snr}) και καταλήγουν σε Terrestrial-only μόνο
όταν ο δορυφόρος βγαίνει εκτός ορατότητας στο τέλος -- μια \emph{αφαίρεση
δευτερεύοντος κόμβου} (SN event), όχι handover: ο σταθμός βάσης που τους
εξυπηρετεί δεν αλλάζει ποτέ. Οι χρήστες $5$-$6$ (πολύ μακριά από κάθε
σταθμό βάσης, BS SNR πάντα βαθιά κάτω από $\text{SNR}_{\min}$) χάνουν την
ίδια ορατότητα δορυφόρου προς το τέλος της διέλευσης, αλλά, δίχως ποτέ
κανέναν χρησιμοποιήσιμο επίγειο υποψήφιο, καταλήγουν σε ρητή κατάσταση
outage αντί σε Terrestrial-only (βλ. Ενότητα~\ref{subsec:results-outage}).
Αυτή η διαφοροποίηση -- ίδιο γεγονός στον δορυφόρο (απώλεια ορατότητας),
διαφορετική έκβαση ανάλογα με το αν υπάρχει χρησιμοποιήσιμος επίγειος
υποψήφιος -- είναι η αναμενόμενη, φυσικά ερμηνεύσιμη συμπεριφορά ενός
ρεαλιστικού κανόνα επιλογής κόμβου. Η σύγκριση με το αρχικό, ασυσχέτιστο
μοντέλο σκίασης (όπου οι ίδιοι χρήστες εμφάνιζαν έως και $29$ handovers
μέσα σε $460$~s) επιβεβαιώνει ότι το τεχνητά συχνό ping-pong της
Ενότητας~\ref{sec:meth-temporal} ήταν πράγματι μεθοδολογικό τεχνούργημα
(artifact) του ανεξάρτητου δειγματισμού σκίασης και όχι φυσικό φαινόμενο,
και ότι διορθώνεται πλήρως με τη χωρική συσχέτιση κατά
Gudmundson~\cite{gudmundson1991correlation}.

\subsection{Ρητή κατάσταση outage}
\label{subsec:results-outage}

Ένα δεύτερο εύρημα εντοπίστηκε αρχικά στο τέλος της διέλευσης, όταν η
γωνία ανύψωσης του δορυφόρου έπεφτε κάτω από το όριο ορατότητας για τους
απομακρυσμένους χρήστες. Στα βήματα $92$-$93$ ($t = 455$-$460$~s), ο
χρήστης $5$, σε απόσταση $\sim\!107$~km από τον πλησιέστερο σταθμό
βάσης, με γωνία ανύψωσης δορυφόρου κάτω από $10^\circ$ (οριακά εκτός
ορατότητας), αναγκαζόταν από τον αρχικό κανόνα -- ένα καθαρό $\argmax$
χωρίς κατώφλι ελάχιστης αποδεκτής ποιότητας -- να ανατεθεί στον
"καλύτερο" επίγειο υποψήφιο, ο οποίος όμως είχε:
\[
\text{PathLoss} = 754.6\text{~dB}, \qquad \text{SNR} = -607.2\text{~dB},
\qquad \text{Capacity} \to 0, \qquad \text{Energy/bit} \to \infty.
\]
Το να "εξυπηρετείται" επίσημα ένας χρήστης από έναν κόμβο με μηδενική
πρακτική χρησιμότητα δεν ήταν σφάλμα υλοποίησης, αλλά άμεση συνέπεια της
απουσίας κατωφλίου στον αρχικό κανόνα επιλογής: επέλεγε πάντα τον
"λιγότερο κακό" υποψήφιο, ακόμα κι όταν κανένας δεν ήταν στην
πραγματικότητα χρησιμοποιήσιμος. Κάτι τέτοιο δεν θα μπορούσε να
παρατηρηθεί στο στατικό σενάριο αναφοράς (Ενότητα~\ref{sec:results-static}),
όπου η γεωμετρία επιλέχθηκε ώστε ο δορυφόρος να παραμένει πάντα ορατός με
άνεση περιθωρίου -- εδώ φάνηκε η αξία της χρονικής προσομοίωσης ως
εργαλείου εντοπισμού ορίων του μοντέλου.

Το εύρημα αυτό αντιμετωπίστηκε ήδη εντός του Μέρους~1, με την προσθήκη
του κατωφλίου $\text{SNR}_{\min} \approx -7.53$~dB στον κανόνα επιλογής
κόμβου (εξίσωση~\eqref{eq:node-selection}, Ενότητα~\ref{sec:meth-node-selection}).
Ξανατρέχοντας την ίδια ακριβώς διέλευση με τον διορθωμένο κανόνα, ο
χρήστης $5$ (και ομοίως ο χρήστης $6$) καταγράφεται πλέον σωστά ως
\texttt{ServingNode="None"} / \texttt{ServingType="Outage"} στα βήματα
$92$-$93$ (Πίνακας~\ref{tab:handovers}) αντί να ανατεθεί στο ανύπαρκτο
BS2: η χωρητικότητα γίνεται ρητά $0$ (όχι απλώς αριθμητικά αμελητέα λόγω
του σχεδόν μηδενικού SNR) και δεν αποδίδεται καθόλου κατανάλωση ισχύος σε
έναν κόμβο που στην πραγματικότητα δεν εξυπηρετεί κανέναν. Τα διαγνωστικά
του "καλύτερου" (αλλά μη-χρησιμοποιήσιμου) υποψηφίου -- ίδιο
$\text{PathLoss} = 754.6$~dB, $\text{SNR} = -607.2$~dB όπως παραπάνω --
παραμένουν διαθέσιμα στην έξοδο για ανάλυση, απλά δεν καθορίζουν πλέον
μια ψευδή ανάθεση εξυπηρέτησης. Σημειώνεται ότι ο χρήστης $4$, σε
αντίστοιχη απόσταση από τον δορυφόρο, \emph{δεν} καταλήγει σε outage
(Ενότητα~\ref{subsec:results-handovers}) -- έχει έναν επίγειο υποψήφιο
αρκετά καλό ώστε να ξεπερνά το $\text{SNR}_{\min}$ σε όλη τη διάρκεια
της διέλευσης, οπότε παραμένει απλώς εξυπηρετούμενος από τον σταθμό
βάσης (μια αφαίρεση δευτερεύοντος κόμβου, όχι απώλεια κάλυψης) όταν ο
δορυφόρος γίνεται αόρατος. Η διαφοροποίηση αυτή είναι ακριβώς το
ζητούμενο: το μοντέλο πλέον διακρίνει "υπάρχει κάλυψη έστω και από μία
μόνο ζεύξη" από "δεν υπάρχει καθόλου κάλυψη", αντί να τα συγχέει.

\section{Αποτελέσματα πρώτου μοντέλου μηχανικής μάθησης}
\label{sec:results-ml}

Πάνω σε ένα σύνολο δεδομένων $200$ τυχαιοποιημένων σεναρίων
($1313$ γραμμές χρηστών, \texttt{rng(42)}) που παρήγαγε ο
\texttt{monteCarloDriver.m} (Ενότητα~\ref{sec:meth-montecarlo}, το
τεκμηριωμένο default του driver), τα δύο μοντέλα της
Ενότητας~\ref{sec:meth-ml} πέτυχαν τις επιδόσεις του
Πίνακα~\ref{tab:ml-results} (grouped split 961/352 γραμμών σε
150/50 σενάρια εκπαίδευσης/ελέγχου). Η κατανομή κλάσεων στο σύνολο
δεδομένων είναι $0\%$ Terrestrial-only, $32\%$ Satellite, $68\%$
DualConnectivity -- η κατηγορία Terrestrial-only \emph{δεν εμφανίζεται
καθόλου}, αφού η ίδια η δορυφορική γεωμετρία/κάλυψη του σεναρίου κάνει
τον δορυφόρο σχεδόν πάντα ταυτόχρονα χρησιμοποιήσιμο όποτε είναι
χρησιμοποιήσιμος και ο επίγειος υποψήφιος (βλ. συζήτηση στο τέλος αυτής
της ενότητας). Οι μετρικές F1/ROC-AUC υπολογίζονται macro-averaged πάνω
στις κλάσεις που όντως εμφανίζονται στο test set (Terrestrial-only
εξαιρείται αυτόματα, αφού έχει μηδενική υποστήριξη παντού).

\begin{table}[htbp]
\centering
\caption{Επιδόσεις πρόβλεψης \texttt{ServingType} (Terrestrial/
Satellite/DualConnectivity), grouped train/test split ανά
\texttt{ScenarioID} (75/25).}
\label{tab:ml-results}
\begin{tabular}{@{}l c c c@{}}
\toprule
Μοντέλο & Accuracy & F1 (macro) & ROC-AUC (macro, OvR) \\
\midrule
Λογιστική Παλινδρόμηση & 0.9858 & 0.9840 & 0.9993 \\
Random Forest          & 0.9972 & 0.9968 & 1.0000 \\
\bottomrule
\end{tabular}
\end{table}

\begin{figure}[htbp]
\centering
\includegraphics[width=\textwidth]{roc_curve.png}
\caption{One-vs-rest καμπύλες ROC για τα δύο μοντέλα, ανά κλάση. Το
υποπλαίσιο Terrestrial δεν έχει καμπύλες, αφού η κλάση αυτή δεν
εμφανίζεται στο test set.}
\label{fig:ml-roc}
\end{figure}

\begin{figure}[htbp]
\centering
\includegraphics[width=0.7\textwidth]{feature_importance_RandomForest.png}
\caption{Σπουδαιότητα χαρακτηριστικών (Random Forest). Τα δύο υποψήφια SNR
κυριαρχούν, όπως αναμενόταν από τον κανόνα κατωφλίου της
εξίσωσης~\eqref{eq:node-selection}.}
\label{fig:ml-feature-importance}
\end{figure}

Το Random Forest παραμένει, όπως προβλέφθηκε στην
Ενότητα~\ref{sec:meth-ml}, ουσιαστικά τέλειο ($99.72\%$ Accuracy,
$100\%$ ROC-AUC): αναμενόμενη συνέπεια του ότι η ίδια η ετικέτα ορίζεται
σχεδόν ντετερμινιστικά από δύο από τα χαρακτηριστικά εισόδου
(\texttt{CandBS\_SNR\_dB}, \texttt{CandSat\_SNR\_dB}), κάτι που
επιβεβαιώνεται και από το Σχήμα~\ref{fig:ml-feature-importance}, όπου τα
δύο αυτά χαρακτηριστικά κυριαρχούν σαφώς έναντι όλων των υπολοίπων. Η
Λογιστική Παλινδρόμηση πλησιάζει πλέον πολύ κοντά ($98.58\%$ Accuracy) --
σε αντίθεση με το αρχικό, δυαδικό πείραμα πριν την πολυσυνδεσιμότητα,
όπου η Λογιστική Παλινδρόμηση υστερούσε αισθητά ($88.35\%$): με τρεις
κλάσεις αντί για δύο, και με την DualConnectivity κλάση να είναι πλέον η
πλειοψηφική ($68\%$), το πρόβλημα ταξινόμησης είναι σε γενικές γραμμές
ευκολότερο για ένα γραμμικό μοντέλο, αφού η κυρίαρχη κλάση είναι και η
πιο εύκολα διαχωρίσιμη (ξεπερνά και τα δύο κατώφλια SNR ταυτόχρονα). Η
αξία αυτού του πρώτου πειράματος βρίσκεται συνεπώς αμιγώς στην
επαλήθευση ορθότητας του συνολικού pipeline δεδομένων
(Ενότητα~\ref{sec:impl-dataflow}), όχι σε απόδειξη ότι το πρόβλημα
πρόβλεψης connectivity είναι εύκολο -- βλ. την επόμενη ενότητα για το πώς
αυτό ερμηνεύεται μαζί με τα άλλα δύο πειράματα.

\section{Προς ρεαλιστικότερα σενάρια πρόβλεψης}
\label{sec:results-ml-realistic}

Το πείραμα της προηγούμενης ενότητας δίνει στο μοντέλο κάτι που ένα
πραγματικό σύστημα δεν θα είχε ποτέ: το ακριβές, στιγμιαίο SNR ενός
υποψήφιου κόμβου με τον οποίο ο χρήστης δεν έχει καν συνδεθεί ακόμα. Για
να εκτιμηθεί πόσο από την υψηλή απόδοση του Πίνακα~\ref{tab:ml-results}
οφείλεται σε αυτή την "κλειδωμένη" πληροφορία, εκπαιδεύτηκαν δύο ακόμα
παραλλαγές των ίδιων δύο ταξινομητών, πάνω στο ίδιο σύνολο δεδομένων και
grouped split, με προοδευτικά λιγότερη πληροφορία SNR. Οι τρεις
παραλλαγές δεν υπάρχουν απλώς για σύγκριση μεταξύ τους -- έχουν η καθεμία
συγκεκριμένο ρόλο: το ακριβές SNR είναι \emph{έλεγχος ορθότητας} του
pipeline (Ενότητα~\ref{sec:bg-ml}), όχι πρόβλεψη προς αναφορά· η
\textbf{μόνο-γεωμετρία} παραλλαγή είναι ένα \emph{κατώτατο όριο} (τι θα
πετύχαινε ένα μοντέλο χωρίς καμία μέτρηση σήματος)· και η
\textbf{θορυβώδης-SNR} παραλλαγή είναι αυτή που η παρούσα εργασία θεωρεί
\emph{αντιπροσωπευτική μιας πραγματικής ανάπτυξης}, αφού προσεγγίζει τι θα
είχε στη διάθεσή του ένα πραγματικό σύστημα (μέτρηση, όχι τέλεια γνώση,
αλλά ούτε πλήρης απουσία σήματος) -- και είναι το αποτέλεσμα στο οποίο
βασίζονται τα συμπεράσματα των Κεφαλαίων~\ref{ch:discussion}
και~\ref{ch:conclusions} όπου αναφέρεται απόδοση πρόβλεψης.

\begin{itemize}
  \item \textbf{Θορυβώδες SNR.} Τα \texttt{CandBS\_SNR\_dB}/
    \texttt{CandSat\_SNR\_dB} παραμένουν χαρακτηριστικά, αλλά πριν
    δοθούν στο μοντέλο προστίθεται σε καθένα ανεξάρτητος γκαουσιανός
    θόρυβος, ώστε να αναπαρασταθεί μια ελαφρώς παλαιωμένη/ατελής μέτρηση
    αντί της ακριβούς τρέχουσας τιμής -- πιο κοντά σε ό,τι θα είχε στη
    διάθεσή του ένα πραγματικό UE μέσω περιοδικών αναφορών μέτρησης
    γειτονικών κυψελών (π.χ. Event A3/A5, TS~38.331) παρά η τέλεια,
    στιγμιαία τιμή του προηγούμενου πειράματος. Η τυπική απόκλιση του
    θορύβου δεν επινοήθηκε: είναι η \emph{πραγματική, μετρημένη}
    διακύμανση SNR ανά ζεύξη από τα $500$ επαναλαμβανόμενα runs της
    Ενότητας~\ref{sec:results-repeated} (Πίνακας~\ref{tab:link-snr-std})
    -- $19.6721$~dB για το επίγειο σκέλος, $0.0682$~dB για το δορυφορικό
    (μικρή τιμή επειδή το δορυφορικό κανάλι εδώ είναι ντετερμινιστικό
    δεδομένης της γεωμετρίας -- βλ. Ενότητα~\ref{sec:meth-sat-channel}).
  \item \textbf{Μόνο γεωμετρία.} Αφαιρούνται εντελώς τα δύο SNR
    χαρακτηριστικά, \emph{και} τα δύο χαρακτηριστικά path loss
    (\texttt{CandBS\_PathLoss\_dB}/\texttt{CandSat\_PathLoss\_dB}), αφού
    το path loss είναι σχεδόν affine συνάρτηση του SNR δεδομένου
    σταθερού EIRP/θορύβου ανά τύπο κόμβου (εξισώσεις~\eqref{eq:snr-bs},
    \eqref{eq:snr-sat}) -- η διατήρησή του θα ξανάφερνε το ίδιο
    "shortcut" υπό άλλο όνομα. Μένουν μόνο \texttt{CandBS\_Distance\_m},
    \texttt{CandSat\_Elevation\_deg}, \texttt{CandSat\_SlantRange\_m},
    \texttt{CandSat\_Visible}, και το context (\texttt{NumBS},
    \texttt{NumUsers}, \texttt{ScenarioType}) -- ό,τι θα ήταν γνωστό σε
    ένα σύστημα \emph{πριν} από οποιαδήποτε μέτρηση σήματος.
\end{itemize}

\begin{table}[htbp]
\centering
\caption{Επιδόσεις πρόβλεψης \texttt{ServingType} με προοδευτικά λιγότερη
πληροφορία SNR (ίδιο dataset/split με τον Πίνακα~\ref{tab:ml-results}).}
\label{tab:ml-results-realistic}
\small
\begin{tabular}{@{}l c c c c@{}}
\toprule
Χαρακτηριστικά εισόδου & LR Accuracy & LR ROC-AUC & RF Accuracy & RF ROC-AUC \\
\midrule
Ακριβές SNR (Πίνακας~\ref{tab:ml-results})  & 0.9858 & 0.9993 & 0.9972 & 1.0000 \\
Θορυβώδες SNR                                & 0.9858 & 0.9973 & 0.9915 & 0.9965 \\
Μόνο γεωμετρία (χωρίς SNR)                   & 0.9830 & 0.9859 & 0.9915 & 0.9902 \\
\bottomrule
\end{tabular}
\end{table}

\begin{figure}[htbp]
\centering
\includegraphics[width=0.7\textwidth]{feature_importance_RandomForest_geometry_only.png}
\caption{Σπουδαιότητα χαρακτηριστικών (Random Forest, μόνο γεωμετρία --
χωρίς SNR ή path loss). Η απόσταση προς το καλύτερο BS κυριαρχεί, αφού
είναι πλέον το χαρακτηριστικό που καθορίζει σχεδόν αποκλειστικά αν ένας
χρήστης είναι Satellite-only ή DualConnectivity -- βλ. συζήτηση παρακάτω.}
\label{fig:ml-feature-importance-geo}
\end{figure}

Το πρώτο πράγμα που ξεχωρίζει στον Πίνακα~\ref{tab:ml-results-realistic}
είναι αυτό \emph{που δεν συμβαίνει}: δεν υπάρχει πια η καθαρή, μονότονη
πτώση Accuracy που θα περίμενε κανείς καθώς αφαιρείται πληροφορία SNR.
Το Random Forest πιάνει $99.15\%$ και στα δύο πειράματα μειωμένης
πληροφορίας -- \emph{το ίδιο} ποσοστό στο θορυβώδες SNR όσο και στη μόνο
γεωμετρία -- ενώ ακόμα και η Λογιστική Παλινδρόμηση μένει πάνω από
$98\%$ και στα τρία πειράματα. Πριν την πολυσυνδεσιμότητα, το ίδιο αυτό
πείραμα έδειχνε μια σαφή, τρίβημη σκάλα ($100\% \rightarrow 89.5\%
\rightarrow 87.2\%$ Accuracy για το Random Forest): η αιτία της
διαφοράς δεν είναι σφάλμα στην υλοποίηση της αξιολόγησης (η μάσκα σε
κλάσεις παρούσες στο test set, Ενότητα~\ref{sec:meth-ml}, ελέγχθηκε
ρητά ώστε να μην κρύβει τεχνητά μια πτώση), αλλά μια πραγματική συνέπεια
της αλλαγής της ίδιας της ετικέτας: αφού η κατηγορία Terrestrial-only
εξαφανίστηκε από το dataset (Ενότητα~\ref{sec:results-ml}), το πρόβλημα
τριών κλάσεων ανάγεται ουσιαστικά σε \emph{DualConnectivity έναντι
Satellite-only} -- και αυτή η διάκριση εξαρτάται σχεδόν αποκλειστικά από
το αν ο χρήστης τοποθετήθηκε "κοντά" ή "μακριά" από σταθμό βάσης στη
γεννήτρια Monte Carlo (Ενότητα~\ref{sec:meth-montecarlo}: πιθανότητα
$30\%$ για "μακρινό" χρήστη, με ακτίνα $1$-$150$~km), \emph{όχι} από μια
λεπτή διαφορά SNR μεταξύ δύο κοντινών υποψηφίων. Η απόσταση BS από μόνη
της (Σχήμα~\ref{fig:ml-feature-importance-geo}) κωδικοποιεί σχεδόν όλη
αυτή την πληροφορία, οπότε η γεωμετρία φτάνει σχεδόν εξίσου μακριά με το
ακριβές SNR. Πρόκειται για ένα πραγματικό εύρημα σχετικά με \emph{αυτή}
τη γεωμετρία ανάπτυξης (ένα σύμπλεγμα σταθμών βάσης που καλύπτεται πλήρως
από το αποτύπωμα ενός δορυφόρου LEO στο όριο ανύψωσης $10^\circ$), όχι
απόδειξη ότι το πρόβλημα connectivity είναι γενικά εύκολο, ούτε
οπισθοδρόμηση της πειραματικής σχεδίασης -- συζητείται περαιτέρω στο
Κεφάλαιο~\ref{ch:discussion}.

Το θορυβώδες SNR εξακολουθεί να προσφέρει κάτι μετρήσιμο, μόνο σε πιο
λεπτό επίπεδο απ' ό,τι πριν: το ROC-AUC του ($0.9965$ RF, $0.9973$ LR)
παραμένει καλύτερο από αυτό της μόνο-γεωμετρίας ($0.9902$ RF, $0.9859$
LR), παρότι οι δύο Accuracy τιμές τους συμπίπτουν ή σχεδόν συμπίπτουν --
η ποιότητα κατάταξης (ranking) των πιθανοτήτων του μοντέλου βελτιώνεται
ακόμα κι όταν το κατώφλι $0.5$ καταλήγει στην ίδια τελική απόφαση. Αυτό
είναι το μόνο σημείο όπου η αναμενόμενη διαβάθμιση ρεαλισμού (ακριβές
$>$ θορυβώδες $>$ γεωμετρία) επιβιώνει καθαρά μετά την πολυσυνδεσιμότητα.

Η διαμόρφωση ενός ακόμα δυσκολότερου, μη-τετριμμένου στόχου πρόβλεψης
(π.χ. πρόβλεψη handover πριν συμβεί, regression/ranking στόχοι, ή ένα
σενάριο γεννήτριας που παράγει πραγματικές Terrestrial-only περιπτώσεις
ώστε το πρόβλημα να μην κυριαρχείται από τη γεωμετρία τοποθέτησης
χρηστών) συζητείται στο Κεφάλαιο~\ref{ch:discussion}.
